% Author: Prof. Dr. Matthias Jung, DL9MJ
% Year: 2026
% Ring modulator, case 2: U2 < 0 and U3 = -U1
\begin{circuitikz}[european, straight voltages]
    \ctikzset{
        quadpoles/transformer core/inner=1,
        quadpoles/transformer core/width=0.6,
        quadpoles/transformer core/height=2.0,
        transformer L1/.style={inductors/width=1.5, inductors/coils=6},
        transformer L2/.style={inductors/width=1.5, inductors/coils=6},
    }

    \begin{scope}[scale=1.25, transform shape, local bounding box=bb]
        \ctikzset{
            diodes/scale=\getDarcImageFactor,
            inductors/scale=\getDarcImageFactor,
            blocks/scale=\getDarcImageFactor
        }
        % T1:
        \draw(-0.5,0) node [transformer core, american](T1){T1};
        \draw(T1.A1)
            to [short,-o] ++(-0.5,0) coordinate(h)
            to [open, v={${\color{DARCblue}U_1}$}] (h|-T1.A2)
            to [short,o-] (T1.A2);

        % Help coordinates:
        \draw(T1.B1)
            to [short,-*] ++(2.5,0) coordinate(north);
        \draw(T1.B2)
            to [short,-*] ++(2.5,0) coordinate(south);
        \draw(north)
            to [open, name={center}] (south);
        \draw(center.center)
            to [open] ++(-1.65,0) coordinate(west);
        \draw(center.center)
            to [open] ++(+1.65,0) coordinate(east);

        % Diode ring:
        \draw(north) to [stroke diode, invert, -*] (west);
        \draw(north) to [stroke diode, -*] (east);
        \draw(south) to [stroke diode, -*] (west);
        \draw(south) to [stroke diode, invert, -*] (east);

        % T2:
        \draw(east)
            to [short] (east|-T1.B2)
            to [short] ++(0.9,0)
            node [transformer core, american, anchor=A2](T2){T2};
        \draw(T2.B1)
            to [short,-o] ++(+0.5,0) coordinate(h)
            to [open, v^={\rotatebox{90}{${\color{DARCblue}\scriptstyle U_3=-U_1}$}}] (h|-T2.B2)
            to [short,o-] (T2.B2);
        \draw(T2.A1)
            to [short] ++(0,0.5) coordinate(h)
            to [short] (h-|west)
            to [short] (west);

        % Local oscillator:
        \draw(T1-L2.midtap)
            to [short, *-] ++(+0.25,0)
            to [short] ++(0,-2.25) coordinate(start);
        \draw(T2-L1.midtap)
            to [short, *-] ++(-0.25,0)
            to [short] ++(0,-2.25) coordinate(end);
        \draw(start)
            to [cgenerator, v={${\color{DARCblue}\scriptstyle U_2=-1\,\mathrm{V}}$}] (end);

        % Node voltages for this switching state:
        \tikzset{node voltage/.style={
            text=DARCblue,
            font=\scriptsize,
            inner sep=1pt
        }}
        \node[node voltage, anchor=north west]
            at ($(T1.B1)+(0.95,0.4)$) {$-1\,\mathrm{V}+\frac{U_1}{2}$};
        \node[node voltage, anchor=north west]
            at ($(T1.B2)+(0.55,-0.2)$) {$-1\,\mathrm{V}-\frac{U_1}{2}$};
        \node[node voltage, anchor=north east]
            at ($(T2.A1)+(-0.35,0.4)$) {$-1\,\mathrm{V}+\frac{U_1}{2}$};
        \node[node voltage, anchor=north east]
            at ($(T2.A2)+(-0.7,-0.2)$) {$-1\,\mathrm{V}-\frac{U_1}{2}$};
        \node[node voltage, anchor=south west]
            at ($(T1-L2.midtap)+(0.18,0.18)$) {$-1\,\mathrm{V}$};
        \node[node voltage, anchor=south east]
            at ($(T2-L1.midtap)+(-0.18,0.18)$) {$0\,\mathrm{V}$};

        % Conducting diode pair and signal paths:
        \draw[DARCred] (T1.B1) to [short] (north);
        \draw[DARCred] (north)
            to [stroke diode, invert, draw=DARCred, fill=DARCred!10, *-*] (west);
        \draw[DARCred] (T2.A1)
            to [short] ++(0,0.5) coordinate(activeTop)
            to [short] (activeTop-|west)
            to [short] (west);
        \draw[DARCred] (T1.B2) to [short] (south);
        \draw[DARCred] (south)
            to [stroke diode, invert, draw=DARCred, fill=DARCred!10, *-*] (east);
        \draw[DARCred] (east) |- (T2.A2);
        % Stable canvas for all three applet states:
        \path (-1.9,-3.5) rectangle (6.92,1.9);
    \end{scope}

    \begin{scope}[shift={($(bb.south west)+(0.75/\getDarcImageFactor,-2.1/\getDarcImageFactor)$)}]
        \begin{axis}[
            axis lines=center,
            axis line style={-Triangle},
            scale only axis,
            axis on top,
            height=1.8cm,
            width=\linewidth*\getDarcImageFactor,
            xmin=0,
            xmax=2.5*pi,
            ymin=-1,
            ymax=1,
            ticks=none,
            xlabel={$t$},
            xlabel style={anchor=north},
            ylabel={$U_2$},
            ylabel style={anchor=east},
        ]
        \pgfplotsinvokeforeach{0,...,11}{
            \fill[DARCblue!10]
                (axis cs:{(2*#1+1)*pi/10},-1)
                rectangle (axis cs:{(2*#1+2)*pi/10},1);
        }
        \addplot[mark=none, DARCblue, thick, domain=0:2.5*pi, samples=1201]
            {0.9*sign(sin(deg(10*x)))};
        \end{axis}
    \end{scope}

    \begin{scope}[shift={($(bb.south west)+(0.75/\getDarcImageFactor,-4.2/\getDarcImageFactor)$)}]
        \begin{axis}[
            axis lines=center,
            axis line style={-Triangle},
            scale only axis,
            axis on top,
            height=1.8cm,
            width=\linewidth*\getDarcImageFactor,
            xmin=0,
            xmax=2.5*pi,
            ymin=-1,
            ymax=1,
            ticks=none,
            xlabel={$t$},
            xlabel style={anchor=north},
            ylabel={$U_1$},
            ylabel style={anchor=east},
        ]
        \pgfplotsinvokeforeach{0,...,11}{
            \fill[DARCblue!10]
                (axis cs:{(2*#1+1)*pi/10},-1)
                rectangle (axis cs:{(2*#1+2)*pi/10},1);
        }
        \addplot[smooth, mark=none, DARCblue, thick, domain=0:2.5*pi, samples=400]
            {0.9*sin(deg(x))};
        \end{axis}
    \end{scope}

    \begin{scope}[shift={($(bb.south west)+(0.75/\getDarcImageFactor,-6.3/\getDarcImageFactor)$)}]
        \begin{axis}[
            axis lines=center,
            axis line style={-Triangle},
            scale only axis,
            axis on top,
            height=1.8cm,
            width=\linewidth*\getDarcImageFactor,
            xmin=0,
            xmax=2.5*pi,
            ymin=-1,
            ymax=1,
            ticks=none,
            xlabel={$t$},
            xlabel style={anchor=north},
            ylabel={$U_3$},
            ylabel style={anchor=east},
        ]
        \pgfplotsinvokeforeach{0,...,11}{
            \fill[DARCblue!10]
                (axis cs:{(2*#1+1)*pi/10},-1)
                rectangle (axis cs:{(2*#1+2)*pi/10},1);
        }
        \addplot[mark=none, DARCblue, thick, domain=0:2.5*pi, samples=1201]
            {0.9*sign(sin(deg(10*x)))*sin(deg(x))};
        \end{axis}
    \end{scope}
\end{circuitikz}
